\section{Data}\label{sec:data}
\paragraph{NetFlow datasets.} We use the v3 release of NF-UNSW-NB15, NF-ToN-IoT, NF-BoT-IoT and NF-CSE-CIC-IDS2018,
which all share one feature set. We drop \res{cfg.data.columns.leaky.count} identifier columns that would let a model
memorise hosts or time instead of behaviour: the IP addresses, the source port, the flow timestamps and the DNS query
id. The destination port is kept. This leaves \res{cfg.data.columns.features.count} features.

From each dataset we draw a stratified sample of \res{cfg.data.sampling.sample_size} flows. Classes with fewer than
\res{cfg.data.sampling.rare_class_threshold} flows are kept whole. NF-BoT-IoT is almost entirely attack traffic, so it
keeps all of its benign flows. After the identifier columns are gone, we remove exact duplicates and flows that have
identical features but different labels. Each sample is then split
\respercentof{cfg.data.split.train}/\respercentof{cfg.data.split.val}/\respercentof{cfg.data.split.test} into train,
validation and test sets, stratified by class (sizes in \cref{tab:data}, \cref{app:data}).

\paragraph{Edge-IIoTset.} We use the published ``DNN'' file and drop the same
\res{cfg.data_edge.columns.leaky.count} columns that the dataset authors drop (addresses, ports, payloads and
timestamps). This leaves \res{cfg.data_edge.columns.features.count} features. Sampling, deduplication and splitting
follow the NetFlow procedure.


\paragraph{Leakage check.} Before modelling, we train XGBoost on each feature alone (train $\to$ validation) and
look closely at any feature that separates the classes by itself. In NF-UNSW-NB15, the minimum TTL alone reaches
binary macro-F1 \res{leakage.nf_unsw_nb15.binary.MIN_TTL}, and the maximum TTL reaches
\res{leakage.nf_unsw_nb15.binary.MAX_TTL}. The reason is that all attack traffic comes from one subnet and all benign
traffic from another, so the TTL and packet-size fingerprints of these machines give the label away. In the
multiclass task the same feature reaches only \res{leakage.nf_unsw_nb15.multiclass.MIN_TTL}, because all attack
classes share the same attacker hosts. We keep all features and treat the binary within-dataset results on
NF-UNSW-NB15 as saturated.
